
\documentclass[conference]{IEEEtran}

\usepackage[T1]{fontenc}
\usepackage[utf8]{inputenc}
\usepackage{graphicx}
\usepackage{amsmath,amssymb}
\usepackage{booktabs}
\usepackage{array}
\usepackage{multirow}
\usepackage{cite}
\usepackage{url}
\usepackage{float}
\usepackage{xcolor}
\usepackage{caption}
\usepackage{setspace}

% Compile-safe figure helpers: if the image file exists, it is inserted;
% otherwise a visible placeholder is shown.
\newcommand{\paperfigure}[5]{%
\begin{figure}[t]
\centering
\IfFileExists{#2}{\includegraphics[width=#1]{#2}}{\fbox{\parbox[c][4.5cm][c]{0.92\columnwidth}{\centering #3}}}%
\caption{#4}
\label{#5}
\end{figure}
}

\newcommand{\paperfigurewide}[5]{%
\begin{figure*}[!t]
\centering
\IfFileExists{#2}{\includegraphics[width=#1]{#2}}{\fbox{\parbox[c][5.8cm][c]{0.92\textwidth}{\centering #3}}}%
\caption{#4}
\label{#5}
\end{figure*}
}

\title{Model-Guided Dataset Expansion for Imbalanced COVID-19 Chest X-Ray Classification}

\author{%
\IEEEauthorblockN{Manas Zambre}
\IEEEauthorblockA{Department of Computer Engineering and Technology \\
MIT World Peace University \\
Pune, India}
}

\begin{document}

\maketitle

\begin{abstract}
Class imbalance remains a persistent issue in medical image classification because deep models are often exposed to many normal findings and comparatively fewer images from less frequent disease categories. This imbalance can produce acceptable aggregate accuracy while leaving weaker performance on clinically important minority classes. This paper studies a model-guided dataset expansion framework for four-class chest X-ray classification on the COVID-19 Radiography Dataset. A pretrained ResNet18 classifier is first trained on real chest X-ray images and then used to identify a weak minority class from validation recall. The selected class is expanded through targeted strategies, including random real-image oversampling and Deep Convolutional Generative Adversarial Network (DCGAN) synthesis, before classifier fine-tuning. The final test comparison includes five settings: baseline training, continued training, focal loss, real-image oversampling, and GAN-based augmentation. Real-image oversampling achieved the best macro F1-score of 0.8806, while GAN augmentation achieved a competitive macro F1-score of 0.8788 and the highest macro recall of 0.897. These results suggest that model-guided expansion can improve class-balanced behavior over the baseline, but they also show that generative augmentation should be evaluated against simple non-generative controls rather than assumed to be superior.
\end{abstract}

\begin{IEEEkeywords}
COVID-19, chest X-ray classification, class imbalance, data-centric deep learning, model-guided augmentation, DCGAN, ResNet18
\end{IEEEkeywords}

\section{Introduction}
Deep learning has become one of the most widely adopted approaches for medical image analysis because convolutional architectures can learn hierarchical spatial patterns directly from image data and often outperform pipelines based on hand-crafted features \cite{lecun2015,litjens2017}. During the COVID-19 pandemic, chest radiography attracted significant research interest because it was inexpensive, widely available, and operationally easier to deploy than laboratory-based testing in many clinical environments. Although RT-PCR remained the reference diagnostic test, chest X-rays were frequently examined for screening support, severity assessment, and rapid triage, particularly when testing resources or reporting speed were limited \cite{chowdhury2020,apostolopoulos2020,ozturk2020}.

One of the most persistent practical problems in medical imaging is class imbalance. Real datasets rarely contain the same number of images for every disease category. Common findings are naturally over-represented, while less frequent categories remain small, expensive to collect, or difficult to annotate. The result is that optimization becomes biased toward the dominant classes, and the model may appear strong overall while still learning weak decision boundaries for the categories that matter most clinically. Standard augmentation such as flips, translations, and mild rotations is still useful for regularization, but it does not necessarily create new pathology-specific variation. In radiographic settings where classes may overlap visually, simply repeating the same appearance under geometric transformations is often not enough \cite{shorten2019,fridadar2018}.

Generative models provide a more interesting alternative because they can synthesize new samples rather than only transform existing ones. Since the original GAN formulation introduced by Goodfellow \textit{et al.} \cite{goodfellow2014}, synthetic data generation has become an active research direction in medical imaging, low-data learning, and imbalance reduction. CovidGAN showed that GAN-generated chest X-ray images could improve COVID-19 detection, but it used a binary formulation and did not let the classifier itself decide which class required additional data \cite{waheed2020}. That distinction matters. In practice, the weakest class is not always the smallest class, and uniform augmentation can spend effort where the classifier is already performing adequately.

This paper therefore adopts a closed-loop, model-guided dataset expansion strategy. A baseline classifier first reveals its weakness through class-wise validation recall. The selected minority class is then expanded through targeted augmentation, and the classifier is fine-tuned while the pretrained ResNet18 backbone remains frozen. Five settings are compared: the original baseline, continued training on the original data, focal loss, oversampling by duplicating real samples, and targeted GAN-based augmentation. The contribution of the work lies less in claiming that a GAN is always superior and more in testing whether classifier feedback can guide where augmentation effort should be applied.

\section{Related Work}
Early COVID-19 chest X-ray studies established that transfer learning was a practical baseline for pandemic-era radiographic classification. Chowdhury \textit{et al.} introduced one of the most widely used COVID-19 chest X-ray datasets and demonstrated that pretrained convolutional models could separate COVID-19, viral pneumonia, lung opacity, and normal classes with encouraging accuracy \cite{chowdhury2020}. Apostolopoulos and Mpesiana reported similarly strong results with transfer-learned CNNs, reinforcing the value of pretrained visual backbones when domain-specific data is limited \cite{apostolopoulos2020}. Ozturk \textit{et al.} proposed DarkCovidNet and further showed that deep neural networks could produce competitive COVID-19 detection performance under constrained data settings \cite{ozturk2020}. These studies defined the classification problem clearly, but they were primarily architecture-oriented and did not focus on targeted dataset expansion.

Larger chest X-ray benchmarks created before COVID-19 also shaped this research area. ChestX-ray8 provided a large hospital-scale dataset for thoracic disease learning and helped normalize deep radiographic benchmarking \cite{wang2017}. CheXNet later showed that deep residual networks could reach very strong pneumonia detection performance and made residual backbones especially influential in chest radiograph analysis \cite{rajpurkar2017}. Those works are relevant here because they support the use of residual transfer learning, but they also indirectly highlight a recurring limitation: strong backbones alone do not solve problems caused by data imbalance and uneven representation.

GAN-based data generation has a much broader history. The original GAN framework introduced adversarial learning as a way to synthesize realistic data \cite{goodfellow2014}. DCGAN improved image generation stability by replacing fully connected adversarial generators with convolutional and transposed-convolutional designs that better preserve spatial structure \cite{radford2015}. Conditional GANs extended the idea by making class-controlled synthesis possible, which is particularly useful when one wants images from a specific disease category \cite{mirza2014}. Progressive GANs and StyleGAN later improved image fidelity even further, although at a higher computational and implementation cost \cite{karras2018,karras2019}. More recently, diffusion models have achieved excellent image quality, but they are heavier and less convenient for a compact course-scale experimental pipeline \cite{ho2020}.

Within medical imaging, several studies have shown that synthetic images can provide real value. Frid-Adar \textit{et al.} reported that GAN-based synthetic augmentation improved liver lesion classification more than conventional augmentation alone \cite{fridadar2018}. Madani \textit{et al.} used GANs in chest radiograph classification and observed that adversarial synthesis could increase training diversity while also helping domain adaptation \cite{madani2018}. Waheed \textit{et al.} proposed CovidGAN specifically for COVID-19 chest X-ray augmentation and showed improved classification in a binary COVID-versus-normal setting \cite{waheed2020}. These studies support the core premise of the present work, namely that synthetic medical images can strengthen learning. At the same time, most prior approaches either augmented a manually chosen class or applied generation without using model performance as the trigger.

Imbalance has also been addressed through non-generative methods. SMOTE and related oversampling strategies synthesize minority samples in feature space and remain influential in classical imbalanced learning, though they do not map naturally onto high-dimensional image structure \cite{chawla2002}. Focal Loss addresses imbalance by shifting optimization emphasis toward harder examples and has become a standard reference point in imbalanced deep classification \cite{lin2017}. Broader augmentation surveys also note that the quality and relevance of new samples matter as much as sample count \cite{shorten2019}. In medical image classification, class weighting, transfer learning, and random oversampling remain common because they are easy to reproduce, even though they may not create enough genuinely new visual variation.

The position of the present work is therefore quite specific. It uses a pretrained residual classifier, a deliberately simple DCGAN generator, a focal-loss comparison, and a feedback step in which the classifier decides which class should be augmented. That last element distinguishes this study from standard augmentation pipelines and from many prior COVID-19 GAN studies where augmentation targets were fixed in advance rather than selected from observed model weakness.

\section{Dataset}
The experiments were conducted on the COVID-19 Radiography Dataset, which contains 21,165 chest X-ray images distributed across four classes: Normal, Lung Opacity, COVID-19, and Viral Pneumonia \cite{chowdhury2020}. The dataset is assembled from multiple institutional and public sources and reflects a naturally imbalanced class distribution rather than an artificially balanced benchmark. The largest class is Normal with 10,192 images, while Viral Pneumonia is the smallest with 1,345 images. This creates a substantial imbalance that is large enough to affect optimization and class-wise generalization.

For the classification pipeline, the dataset was split in a stratified manner into training, validation, and test subsets using an 80/10/10 ratio. Stratification ensured that each split preserved the original class proportions, which is important when evaluating minority-class recall. The resulting split statistics are shown in Table~\ref{tab:dataset}.

\begin{table}[t]
\caption{Dataset statistics and stratified split}
\label{tab:dataset}
\centering
\scriptsize
\setlength{\tabcolsep}{3pt}
\resizebox{0.98\columnwidth}{!}{%
\begin{tabular}{lccccc}
\toprule
Class & Total & Train & Val. & Test & Relative Size \\
\midrule
Normal & 10,192 & 8,153 & 1,019 & 1,019 & 1.0$\times$ \\
Lung Opacity & 6,012 & 4,809 & 601 & 602 & 1.7$\times$ smaller \\
COVID-19 & 3,616 & 2,893 & 361 & 362 & 2.8$\times$ smaller \\
Viral Pneumonia & 1,345 & 1,076 & 135 & 134 & 7.6$\times$ smaller \\
\midrule
Total & 21,165 & 16,932 & 2,116 & 2,117 & --- \\
\bottomrule
\end{tabular}}
\end{table}

\paperfigure{0.96\columnwidth}{fig1_class_distribution.png}{Figure 1 placeholder\\Class distribution bar chart}{Class distribution of the COVID-19 Radiography Dataset.}{fig:class_dist}

\section{Methodology}
\subsection{Pipeline Overview}
The proposed pipeline has three connected stages. First, a baseline classifier is trained on the real dataset and evaluated on the validation split. Second, the weakest minority class is selected using the classifier's own class-wise recall. Third, that selected class is expanded through targeted strategies, including real-image oversampling and DCGAN-generated synthetic samples, before classifier fine-tuning. This gives the method a data-centric feedback structure: the model is not only trained on the data, it also guides how the data should be expanded.

\paperfigurewide{\textwidth}{fig2_architecture.png}{Figure 2 placeholder\\System architecture diagram of the model-guided dataset expansion pipeline}{Overview of the model-guided dataset expansion pipeline.}{fig:pipeline}

\subsection{Baseline Classifier}
A pretrained ResNet18 model was used as the classification backbone because residual networks remain reliable and well understood for transfer learning in medical image tasks \cite{he2016}. The convolutional backbone was frozen and only the classifier head was trained. The head consisted of a fully connected layer mapping 512 features to 256 hidden units, followed by ReLU activation, dropout of 0.4, and a final linear layer producing logits for the four target classes. This configuration kept the number of trainable parameters modest while preserving the benefit of pretrained visual features.

Training was performed for 8 epochs using weighted cross-entropy loss. The class weights were included to partially compensate for imbalance during optimization, but they were not expected to solve the representation problem entirely. Adam was used as the optimizer with learning rate $1\times10^{-3}$ and batch size 32. The baseline model served two purposes: it established the initial classification performance and it produced the recall scores used for weak-class selection.

\subsection{Weak-Class Selection}
The guidance rule was deliberately simple. After baseline training, class-wise recall was computed on the validation split. Instead of automatically choosing the globally smallest class, the method compared the recall values among the minority candidates and selected the one that was hardest for the model to recover consistently. In the observed experiment, COVID-19 achieved lower validation recall than Viral Pneumonia among the small classes, so COVID-19 was selected as the target for synthetic expansion. This decision rule reflects the central idea of the paper: the classifier's own weakness, not just raw class count, should drive augmentation.

\subsection{DCGAN Architecture and Training}
A Deep Convolutional Generative Adversarial Network was trained only on the selected weak class. The generator accepted a 100-dimensional latent noise vector and mapped it through a stack of transposed convolution layers to produce a $64\times64$ RGB image. Batch normalization and ReLU activations were used in intermediate layers, with Tanh activation at the output. The discriminator used standard convolution layers with LeakyReLU activations and a sigmoid output to distinguish real images from synthetic ones \cite{radford2015}.

The GAN was trained for 60 epochs on 2,893 real COVID-19 chest X-ray images. To improve training stability, label smoothing was used with real labels set to 0.9, and the discriminator was updated at a slightly reduced pace relative to the generator. Adam optimization with $\beta_1=0.5$ and $\beta_2=0.999$ was used, following conventional DCGAN practice. A total of 800 synthetic COVID-19 images were generated after training and added to the classifier training set.

\begin{table}[t]
\caption{DCGAN training configuration}
\label{tab:gan}
\centering
\small
\begin{tabular}{lc}
\toprule
Parameter & Value \\
\midrule
Latent dimension & 100 \\
Generator output size & $64\times64$ RGB \\
Training images & 2,893 (COVID-19 class) \\
Epochs & 60 \\
Generator learning rate & $2\times10^{-4}$ \\
Discriminator learning rate & $1\times10^{-4}$ \\
Optimizer & Adam ($\beta_1=0.5$, $\beta_2=0.999$) \\
Synthetic images generated & 800 \\
\bottomrule
\end{tabular}
\end{table}

\paperfigure{0.96\columnwidth}{fig3_real_vs_synthetic.png}{Figure 3 placeholder\\Real versus synthetic COVID-19 chest X-ray samples}{Representative real and GAN-generated COVID-19 chest X-ray images.}{fig:realsynth}

\subsection{Ablation Settings and Fine-Tuning}
To isolate the contribution of targeted expansion more clearly, four post-baseline settings were evaluated. The first was \textit{Continue Only}, in which the baseline classifier was trained for 5 additional epochs on the original data without any new samples. This controlled for the possibility that extra optimization time alone was responsible for improvement. The second was \textit{Focal Loss}, which replaced weighted cross-entropy with focal loss during fine-tuning and served as a stronger imbalance-learning baseline. The third was \textit{Real Oversample}, where 800 real COVID-19 training images were duplicated randomly and added back into training. This controlled for the effect of increased class count without creating new visual content. The fourth was \textit{GAN Augmented}, in which 800 DCGAN-generated COVID-19 images were appended to the original training data.

In all post-baseline settings, the ResNet18 backbone remained frozen and only the classifier head was fine-tuned. This decision kept the experiments lightweight and made the comparison cleaner. If the full backbone had been updated, the interpretation would have mixed the effect of new data with deep feature re-optimization. The same fine-tuning schedule was maintained across all post-baseline settings, with the loss function changed only for the focal-loss baseline.

\section{Results}
\subsection{Overall Performance}
The main quantitative comparison is shown in Table~\ref{tab:main_results}. The baseline model achieved 87.10\% accuracy and a macro F1-score of 0.8724. Continued training improved macro F1 to 0.8773, indicating that the classifier still benefited from a small amount of extra optimization even without new data. Focal loss achieved 87.20\% accuracy and macro F1-score 0.8742. Real oversampling produced the best macro F1-score of 0.8806 and the highest COVID-19 recall of 92.54\%. GAN augmentation achieved a competitive macro F1-score of 0.8788 and the highest macro recall of 0.897.

The gains are modest, but they are meaningful in the context of an imbalanced dataset. Accuracy is not sufficient on its own because it is heavily influenced by the larger Normal and Lung Opacity classes. Macro F1 and macro recall are more informative because they weight the contribution of each class more evenly. The result does not show that GAN augmentation is universally superior to simple oversampling; instead, it shows that model-guided expansion improves over the baseline and that synthetic expansion remains competitive with strong non-generative controls.

\begin{table}[t]
\caption{Main experimental results on the held-out test set}
\label{tab:main_results}
\centering
\scriptsize
\setlength{\tabcolsep}{3pt}
\resizebox{0.98\columnwidth}{!}{%
\begin{tabular}{lcccc}
\toprule
Method & Accuracy & Macro Recall & Macro F1 & COVID Recall \\
\midrule
Baseline & 87.10\% & 0.891 & 0.8724 & 88.40\% \\
Continue Only & 87.58\% & 0.895 & 0.8773 & 91.16\% \\
Focal Loss & 87.20\% & 0.880 & 0.8742 & 90.06\% \\
Real Oversample & \textbf{87.72\%} & 0.895 & \textbf{0.8806} & \textbf{92.54\%} \\
GAN Augmented & 87.67\% & \textbf{0.897} & 0.8788 & 90.33\% \\
\bottomrule
\end{tabular}}
\end{table}

\subsection{Per-Class Analysis}
A more detailed class-wise view of the GAN-augmented setting is shown in Table~\ref{tab:classwise}. The GAN-augmented model achieved precision 0.867, recall 0.903, and F1-score 0.885 on COVID-19. Viral Pneumonia reached recall 0.970 and F1-score 0.900, although it was not directly augmented. Lung Opacity remained comparatively weaker, which is understandable because it overlaps with multiple abnormal radiographic presentations.

An important nuance appears when the ablation settings are compared conceptually. Real oversampling achieved the strongest COVID-19 recall and the highest macro F1, suggesting that repeated exposure to the weak class can be highly effective in this dataset. GAN augmentation did not surpass that result, but it improved over the baseline and achieved the highest macro recall. In simpler terms, the generative strategy was useful, but the experiment also shows why synthetic augmentation must be compared against simple baselines before making strong claims.

\begin{table}[t]
\caption{Per-class results for the GAN-augmented model}
\label{tab:classwise}
\centering
\small
\begin{tabular}{lcccc}
\toprule
Class & Precision & Recall & F1 & Support \\
\midrule
COVID-19 & 0.867 & 0.903 & 0.885 & 362 \\
Lung Opacity & 0.842 & 0.831 & 0.836 & 602 \\
Normal & 0.907 & 0.882 & 0.895 & 1,019 \\
Viral Pneumonia & 0.839 & 0.970 & 0.900 & 134 \\
\midrule
Macro Avg. & 0.864 & 0.897 & 0.879 & 2,117 \\
\bottomrule
\end{tabular}
\end{table}

\paperfigure{0.96\columnwidth}{fig4_confusion_matrices.png}{Figure 4 placeholder\\Confusion-matrix comparison across selected evaluated settings}{Confusion matrices for Baseline, Continue Only, Real Oversample, and GAN Augmented settings.}{fig:confusion}


\section{Discussion}
The most important outcome of the study is that model-guided expansion improved class-balanced performance over the baseline, but the type of expansion mattered. Real-image oversampling produced the best macro F1, while GAN augmentation produced the highest macro recall and remained competitive with the strongest control. This is a more careful result than simply claiming that synthetic images are better. It suggests that the model-guided selection step is useful, but that the augmentation method must be evaluated empirically rather than assumed to be superior because it is generative.

The comparison with real oversampling is especially informative. Random duplication of real COVID-19 images achieved the highest COVID-19 recall, which suggests that repeated minority exposure can strongly increase sensitivity for the selected class. GAN augmentation, however, gave the highest macro recall, indicating that it may distribute recall gains more evenly across classes. This distinction is important for imbalanced medical classification: the best method depends on whether the priority is one selected disease class or overall balance across all classes.

There are still clear limitations. First, the current result is based on a single experimental run, so repeated trials across multiple random seeds are needed before making a strong statistical claim. Second, the class imbalance in the dataset is moderate rather than extreme, so the margin for improvement is naturally limited. Third, the DCGAN outputs are only $64\times64$, which is well below the visual resolution at which subtle radiographic details are usually interpreted. Fourth, the guidance signal is intentionally simple and relies only on class-wise validation recall. A stronger research version of this idea could use sample-level loss, calibration error, uncertainty estimates, or feature-cluster failure analysis to decide where augmentation is needed most.

Even with those limitations, the framework remains useful because it is practical, interpretable, and data-centric. Many small-scale classification studies focus almost entirely on changing the classifier while leaving the training distribution untouched. This study shows that using the classifier to guide where the dataset should be expanded can produce measurable gains without requiring a more complicated backbone. This represents the principal practical contribution of the work.

\section{Conclusion}
This paper presented a model-guided dataset expansion framework for four-class COVID-19 chest X-ray classification. A pretrained ResNet18 classifier was first used to identify the weakest minority class through validation recall. The selected class was then expanded using targeted real-image oversampling and DCGAN-based synthetic augmentation, and the resulting models were compared against continued training and focal loss.

The central idea is straightforward and useful: the classifier should help decide where augmentation effort is spent. In the final test results, real-image oversampling achieved the best macro F1-score of 0.8806, while GAN augmentation remained competitive with a macro F1-score of 0.8788 and achieved the highest macro recall of 0.897. Therefore, the study does not prove that GAN augmentation is universally better than simpler methods. Instead, it supports a more practical conclusion: model-guided expansion can improve class-balanced performance, and generative augmentation is a viable but not automatically dominant expansion strategy.

Future work can extend this direction by running multiple random seeds, adding external validation, moving to higher-resolution generation, testing conditional or diffusion-based synthesis, and using more refined guidance signals based on sample-level or feature-level failure analysis.

\begin{thebibliography}{99}

\bibitem{goodfellow2014}
I. Goodfellow, J. Pouget-Abadie, M. Mirza, B. Xu, D. Warde-Farley, S. Ozair, A. Courville, and Y. Bengio, ``Generative Adversarial Nets,'' in \textit{Advances in Neural Information Processing Systems (NeurIPS)}, 2014, pp. 2672--2680.

\bibitem{mirza2014}
M. Mirza and S. Osindero, ``Conditional Generative Adversarial Nets,'' arXiv:1411.1784, 2014.

\bibitem{radford2015}
A. Radford, L. Metz, and S. Chintala, ``Unsupervised Representation Learning with Deep Convolutional Generative Adversarial Networks,'' arXiv:1511.06434, 2015.

\bibitem{lecun2015}
Y. LeCun, Y. Bengio, and G. Hinton, ``Deep Learning,'' \textit{Nature}, vol. 521, no. 7553, pp. 436--444, 2015.

\bibitem{he2016}
K. He, X. Zhang, S. Ren, and J. Sun, ``Deep Residual Learning for Image Recognition,'' in \textit{Proc. IEEE Conf. Comput. Vis. Pattern Recognit. (CVPR)}, 2016, pp. 770--778.

\bibitem{wang2017}
X. Wang, Y. Peng, L. Lu, Z. Lu, M. Bagheri, and R. M. Summers, ``ChestX-ray8: Hospital-Scale Chest X-Ray Database and Benchmarks on Weakly-Supervised Classification and Localization of Common Thorax Diseases,'' in \textit{Proc. IEEE CVPR}, 2017, pp. 2097--2106.

\bibitem{rajpurkar2017}
P. Rajpurkar, J. Irvin, K. Zhu, \textit{et al.}, ``CheXNet: Radiologist-Level Pneumonia Detection on Chest X-Rays with Deep Learning,'' arXiv:1711.05225, 2017.

\bibitem{lin2017}
T.-Y. Lin, P. Goyal, R. Girshick, K. He, and P. Dollar, ``Focal Loss for Dense Object Detection,'' in \textit{Proc. IEEE ICCV}, 2017, pp. 2980--2988.

\bibitem{litjens2017}
G. Litjens, T. Kooi, B. E. Bejnordi, \textit{et al.}, ``A Survey on Deep Learning in Medical Image Analysis,'' \textit{Medical Image Analysis}, vol. 42, pp. 60--88, 2017.

\bibitem{fridadar2018}
M. Frid-Adar, E. Klang, M. Amitai, J. Goldberger, and H. Greenspan, ``Synthetic Data Augmentation Using GAN for Improved Liver Lesion Classification,'' in \textit{Proc. IEEE ISBI}, 2018, pp. 289--293.

\bibitem{madani2018}
A. Madani, M. Moradi, A. Karargyris, and T. Syeda-Mahmood, ``Semi-Supervised Learning with Generative Adversarial Networks for Chest X-Ray Classification with Ability of Data Domain Adaptation,'' in \textit{Proc. IEEE ISBI}, 2018, pp. 1038--1042.

\bibitem{karras2018}
T. Karras, T. Aila, S. Laine, and J. Lehtinen, ``Progressive Growing of GANs for Improved Quality, Stability, and Variation,'' in \textit{Proc. ICLR}, 2018.

\bibitem{shorten2019}
C. Shorten and T. M. Khoshgoftaar, ``A Survey on Image Data Augmentation for Deep Learning,'' \textit{Journal of Big Data}, vol. 6, no. 60, 2019.

\bibitem{karras2019}
T. Karras, S. Laine, and T. Aila, ``A Style-Based Generator Architecture for Generative Adversarial Networks,'' in \textit{Proc. IEEE CVPR}, 2019, pp. 4401--4410.

\bibitem{chowdhury2020}
M. E. H. Chowdhury, T. Rahman, A. Khandakar, \textit{et al.}, ``Can AI Help in Screening Viral and COVID-19 Pneumonia?,'' \textit{IEEE Access}, vol. 8, pp. 132665--132676, 2020, doi: 10.1109/ACCESS.2020.3021570.

\bibitem{apostolopoulos2020}
I. D. Apostolopoulos and T. A. Mpesiana, ``COVID-19: Automatic Detection from X-Ray Images Utilizing Transfer Learning with Convolutional Neural Networks,'' \textit{Physical and Engineering Sciences in Medicine}, vol. 43, pp. 635--640, 2020.

\bibitem{ozturk2020}
T. Ozturk, M. Talo, E. A. Yildirim, U. B. Baloglu, O. Yildirim, and U. R. Acharya, ``Automated Detection of COVID-19 Cases Using Deep Neural Networks with X-Ray Images,'' \textit{Computers in Biology and Medicine}, vol. 121, 103792, 2020.

\bibitem{waheed2020}
A. Waheed, M. Goyal, D. Gupta, A. Khanna, A. Al-Turjman, and P. R. Pinheiro, ``CovidGAN: Data Augmentation Using Auxiliary Classifier GAN for Improved Covid-19 Detection,'' \textit{IEEE Access}, vol. 8, pp. 91916--91923, 2020, doi: 10.1109/ACCESS.2020.2994762.

\bibitem{ho2020}
J. Ho, A. Jain, and P. Abbeel, ``Denoising Diffusion Probabilistic Models,'' in \textit{Advances in Neural Information Processing Systems (NeurIPS)}, 2020.

\bibitem{chawla2002}
N. V. Chawla, K. W. Bowyer, L. O. Hall, and W. P. Kegelmeyer, ``SMOTE: Synthetic Minority Over-sampling Technique,'' \textit{Journal of Artificial Intelligence Research}, vol. 16, pp. 321--357, 2002.

\end{thebibliography}

\end{document}
